\section{Preliminaries}\label{sec: prelims}
In this section we recall some necessary background material.  
\subsection{\texorpdfstring{Review of $\mathrm{ETHH}$}{Review of ETHH}}

Let $G$ be a finite group. We write $\Spg$ for the category of genuine orthogonal $G$-spectra, and $\Alg(\Spg)$ for the category of associative ring $G$-spectra. We recall the definition of $\mathrm{ETHH}$ from \cite{CGK25}.

\begin{definition}
  Let $R$ be an associative algebra in $\Spg$. The \textit{cyclic bar construction} of $R$, $N^{\cyc}_\bullet(R)$, is a simplicial object given in degree $n$ by $R^{\smsh n+1}$ with face and degeneracy maps as follows:
  \begin{alignat*}{3}
    d_0 &= (\id^{\smsh n-1} \smsh \mu) \circ \tau  \\
    d_i &= (\id^{\smsh i-1}\smsh \mu \smsh \id^{\smsh n-i}) \qquad& 0&< i<n \\
    s_i &= \id^{\smsh i+1} \smsh \eta \smsh \id^{\smsh n-i} \qquad& 0&\leq i<n, 
  \end{alignat*}
  where $\mu$ and $\eta$ are the multiplication and unit in $R$, respectively, and $\tau$ rotates the first factor of $R$ to the end.
\end{definition}

\begin{definition}[{\cite[Definition 3.1]{CGK25}}]
  Let $\mathcal{U}$ be a complete $(G\times S^1)$-universe, and let $\mathcal{V} = i^*_G\mathcal{U}$ be its restriction to a complete $G$-universe.
  Equivariant topological Hochschild homology is the functor $\ETHH\colon\Alg(\Spg) \to \Spg[G\times S^1]$ defined as
  \[ \ETHH(-) = \mathcal{I}^{\mathcal{U}}_{\mathcal{V}}|N^\cyc_\bullet(-)| \]
  where $\mathcal{I}^{\mathcal{U}}_{\mathcal{V}}$ is the change of universe functor (\cite[Section II.1]{LMSM}) and $| \cdot |$ denotes geometric realization.
\end{definition}

\begin{proposition}
    \label{ETHHisCyclicBar}
  Let $R$ be a ring $G$-spectrum. After restricting to a $G$-spectrum, $\ETHH(R)$ agrees with the cyclic bar construction in the category of $G$-spectra:
  \[ \Res^{G\times S^1}_{G\times 1}(R) \we |N^\cyc_\bullet(R)|. \]
\end{proposition}
\begin{proof}
  This follows as restriction is (strong) symmetric monoidal and preserves geometric realization.
\end{proof}

\begin{proposition}
  Equivariant topological Hochschild homology preserves commutative algebras, that is, $G$-$\mathbb{E}_\infty$-algebras.
\end{proposition}
\begin{proof}
  This is \cite[Proposition 3.5]{CGK25}.
\end{proof}

\subsection{Review of equivariant Thom constructions}
To compute $\ETHH$ of the equivariant cobordism spectra $MU_\R$ and $MU_G$ (for $G$ a finite abelian group), we will use the fact that these are equivariant Thom spectra of $G$-$\mathbb{E}_\infty$ maps, and that $\ETHH$ commutes with the equivariant Thom spectrum functor. A point-set model of the equivariant Thom spectra functor was constructed by Lewis and May in X.3.1 of \cite{LMSM}. For our purposes, we will need the equivariant Thom spectrum functor to be $G$-symmetric monoidal, so we will use the model constructed in \cite{HHKWZ24}.

\medskip

We denote the ($G$-category of) $G$-spaces by $\Top _G$, the ($G$-category of) genuine $G$-spectra by $\Spg$, and by $\Pic(\Spg)$ the Picard $G$-space of invertible $G$-spectra. The equivariant Thom spectrum functor takes a map of $G$-spaces $f: X \to \Pic(\Spg)$ and produces a genuine $G$-spectrum $\mathbf{Th}(f)$. We recall:

\begin{theorem}\label{thm-eThom} \cite{HHKWZ24}
    Let $G$ be a finite group. The equivariant Thom spectrum functor 
    $$\mathbf{Th}: \Top _G / \Pic(\Spg) \to \Spg$$ 
    is $G$-symmetric monoidal and commutes with colimits.
\end{theorem}

For us, the main use of this theorem will be in the following examples.

\begin{example}\label{ex-MU}
$\, $

    \begin{enumerate}
        \item Let $G = C_2$. The Real complex cobordism spectrum $MU_\R$ is the equivariant Thom spectrum of the Real J-homomorphism $BU_\R \to \Pic(\Spg)$, which is a $G$-$\mathbb{E}_\infty$ map. (See, for example, Remark 13 of \cite{HL}.)
        \item Let $G$ be any finite group. The equivariant complex cobordism spectrum $MU_G$ is the equivariant Thom spectrum of the equivariant J-homomorphism $BU_G \to \Pic(\Spg)$, which is a $G$-$\mathbb{E}_\infty$ map. (See, for example, Example 6.1.53 of \cite{Sch}.)
    \end{enumerate}

\end{example}

The theorem implies that both $MU_{\R}$ and $MU_G$ have the structure of $G$-commutative monoids in $\Sp_G$. Equivalently, they can be modeled by $G$-$\mathbb{E}_{\infty}$-ring spectra.

\subsection{Review of equivariant factorization homology}
In this subsection, we review the basic properties of equivariant factorization homology that we will need in order to compute $\ETHH(MU_\R)$ and $\ETHH(MU_G)$. 

\medskip

Factorization homology is a homology theory of manifolds, taking in an $n$-manifold $M$ equipped with a trivialization of its tangent bundle and an $\mathbb{E}_n$-algebra $A$ in a category $\mathcal{C}$ and producing $\int_M A \in \mathcal{C}$. Factorization homology is symmetric monoidal (in $M$ under disjoint union, and in $A$ under the symmetric monoidal structure in $\mathcal{C}$), functorial (in open embeddings of $n$-manifolds, and in $\mathbb{E}_n$-algebra maps), and satisfies $\otimes$-excision for collar gluings $M = M' \cup_{M_0 \times \R} M''$. Factorization homology can be defined as a left Kan extension.

\begin{example}
    If $R$ is a ring spectrum, then $\int_{S^1} R \simeq \THH (R)$.
\end{example}


In \cite{Hor}, Horev produces a genuine equivariant version of factorization homology, and shows that various equivariant versions of $\THH$ can be studied using genuine equivariant factorization homology. In this paper, we will show that the restriction of $\ETHH$ to a $G$-spectrum is an instance of equivariant factorization homology, and use this to compute the $\ETHH$ of equivariant complex cobordism spectra. The main input we will need is Theorem 3 of \cite{HHKWZ24}:

\begin{theorem}\label{thm-FH-Thom-HKZ}\cite{HHKWZ24}
    Equivariant Thom spectra commute with equivariant factorization homology. That is, if $f: \Omega ^V X \to \Pic(\Spg)$ is a map of $V$-fold loop spaces (for $V$ a finite-dimensional $G$-representation), then 
    $$\int_M \mathbf{Th}(f) \simeq \mathbf{Th}\left( \int_M \Omega^V X \to \int_M \Pic(\Spg) \to \Pic(\Spg)\right)$$
    for every $G$-manifold $M$ equipped with an equivariant isomorphism of bundles $TM \cong M \times V$.
\end{theorem}

Along with Theorem 4 of \cite{HHKWZ24}, which identifies $\int_M \Omega^V X \simeq \mathrm{Map}(M^+, X)$, we will obtain a multiplicative description of $\int_M$ of Thom spectra of $G$-$\mathbb{E}_\infty$ maps, and apply this to $\ETHH$.

